\chapter{Formal Safeguards with TLA+}
\label{chap:formal-safeguards}
\newcommand{\tlbreak}[2]{\texttt{#1}\allowbreak\texttt{#2}}

CharlotteOS uses executable TLA+ models to review stateful functionality. The
models state each safety contract, explore action orderings in a finite instance,
retain known-bad designs as counterexamples, and map abstract transitions back to
Rust.

This is stronger than relying on selected test schedules, but narrower than a
proof of the complete operating system. TLC proves invariants of the modeled
state machine for the configured finite bounds. The repository supplements
that result with an action-level conformance map, Rust regression tests,
target builds, Clippy, and boot tests. Each layer answers a different question:

\begin{description}
    \item[TLA+ model] Is the protocol design safe under every modeled action
        ordering, crash point, retry, and stale event within the configured
        bounds?
    \item[Conformance map] Which concrete Rust state and linearization point
        implement each abstract action?
    \item[Rust and boot tests] Does the implementation exercise the intended
        transition on the host and on CharlotteOS targets?
    \item[Build and lint gates] Does the integrated implementation remain
        type-correct and warning-free for the actual target?
\end{description}

The models and their configurations live in \repodir|docs/tla|. The concise
model inventory is in \repofile|docs/tla/README.md|; the concrete-to-abstract
mapping is in \repofile|docs/tla/CONFORMANCE.md|. Those files are maintained
with the implementation and are authoritative over the overview in this
chapter.

\section{Why model functionality as a state machine?}

Most failures addressed by these models are not computation errors. They are
ordering errors: a stale cleanup arrives after a replacement, a thread is
removed while it can still execute, a session identifier is reused after a
restart, or a joiner forgets which cluster has authority over it. A unit test
can exercise one such interleaving. A state-machine model asks TLC to enumerate
all reachable interleavings in a deliberately small world.

A TLA+ safety specification has four essential parts:

\begin{enumerate}
    \item \textbf{Variables} record only the state relevant to the property,
        such as an owner, generation, durable term, queue, or admission anchor.
    \item \textbf{Init} defines every allowed initial state.
    \item \textbf{Next} is a disjunction of named actions. Each action defines
        its enabling condition and complete next-state update.
    \item \textbf{Invariants} state what must be true in every reachable state.
\end{enumerate}

The common shape is:

\begin{verbatim}
VARIABLES phase, generation, owner
vars == <<phase, generation, owner>>

Init == ...
Next == Start \/ Replace \/ Stop \/ Crash \/ Restart
Spec == Init /\ [][Next]_vars

TypeOK == ...
Safety == ...
Invariants == TypeOK /\ Safety
\end{verbatim}

The square temporal operator means that every step must be a \texttt{Next}
step or a stuttering step and that the invariant must hold throughout the
behavior. TLC constructs the reachable state graph rather than sampling it.
Constants in a \texttt{.cfg} file bound identities, queue capacities,
generations, terms, and log indices so that this graph is finite.

Small bounds are intentional. Many protocol violations need only two owners,
two generations, or three Raft nodes. Exhaustive exploration of that small
instance is often more revealing than a large randomized run because TLC
systematically covers the orderings rather than hoping to encounter them.

\section{The safeguarding loop}

CharlotteOS applies the models as a feedback loop between design,
implementation, and validation:

\begin{enumerate}
    \item \textbf{Choose a safety boundary.} Identify the authority or resource
        whose invalid transition would be harmful: reply authority, memory
        ownership, an executing address space, a committed log entry, or a
        service generation.
    \item \textbf{Project concrete state.} Remove bytes, addresses, and policy
        details that cannot affect the selected invariant. Preserve identity,
        ownership, durability, ordering, and lifecycle phases.
    \item \textbf{Name the transitions.} Actions use implementation language.
        Names such as \texttt{Receive}, \texttt{Reap}, and \texttt{Restart}
        make review against Rust possible.
    \item \textbf{State the invariant positively.} Examples are ``a reaped
        thread is off-CPU,'' ``a committed entry agrees at every replica,'' and
        ``a stale generation cannot remove its replacement.''
    \item \textbf{Model crashes and stale work.} Durable and volatile state are
        separated where restart matters. Delayed messages, old handles, and
        retries are actions rather than assumptions that the environment is
        well behaved.
    \item \textbf{Retain a negative specification.} When a faulty transition is
        understood, an \texttt{Unsafe...} action and configuration are kept.
        The check suite requires TLC to find the named invariant violation.
        This proves that the invariant and test harness are capable of seeing
        the regression rather than passing vacuously.
    \item \textbf{Map back to code.} \repofile|docs/tla/CONFORMANCE.md| records
        the concrete function, persisted field, and linearization point for
        every modeled action. Disagreement is resolved by changing the model,
        changing CharlotteOS, or documenting behavior outside the abstraction.
    \item \textbf{Validate both levels.} Run TLC, focused Rust tests, target
        builds and Clippy, and relevant boot or multi-node tests.
\end{enumerate}

The model is an executable safety contract, not a diagram added after
implementation. A change is incomplete if
the model claims an atomic or durable transition that the code does not
provide.

\section{What is modeled}

The current suite contains eighteen safe specifications. They are divided by
state ownership and failure boundary so that each model remains reviewable.
The Raft models are deliberately layered: election, log replication,
membership, pre-membership admission, and snapshots have different state and
different proof obligations.

\small
\begin{longtable}{|p{0.26\textwidth}|p{0.26\textwidth}|p{0.38\textwidth}|}
\hline
\textbf{Model} & \textbf{Functional boundary} & \textbf{Principal safety properties} \\
\hline
\endfirsthead
\hline
\textbf{Model} & \textbf{Functional boundary} & \textbf{Principal safety properties} \\
\hline
\endhead
\tlbreak{Charlotte}{IPC} & Endpoint calls, reply tokens, memory move/copy/borrow,
    cancellation, endpoint close, domain teardown & Bounded queues; unique
    receiver-visible reply authority; exclusive write borrow; every remaining
    borrow backed by a live token; no dangling moved-memory capability. \\
\hline
\tlbreak{Charlotte}{EndpointObservers} & Readiness observers and connection-close
    watches & Ordinary message arrival cannot falsely signal endpoint closure;
    closure wakes the appropriate observers. \\
\hline
\tlbreak{Charlotte}{CQ} & Completion operations, CQ ring/backlog, wait/wake,
    cancellation, timers & Non-lossy completion tracking; bounded visible ring;
    no lost wakeup; timer completion remains tied to its operation. \\
\hline
\tlbreak{Charlotte}{Scheduler} & Spawn, admission, execution, blocking, migration,
    cross-LP abort, reaping, address-space destruction & At most one execution
    per LP; only the owner LP retires an executing thread; reaping occurs
    off-CPU; destroyed address spaces have no live threads. \\
\hline
\tlbreak{Charlotte}{AddressSpace} & Reusable numeric ASIDs and software-generation
    handles & A stale handle cannot close a replacement that reused the same
    numeric ASID. \\
\hline
\tlbreak{Charlotte}{HardwareAsid} & Hardware-ASID allocation, retirement, cached
    translations, invalidation & A hardware tag is not reused while stale
    translations for its former owner remain. \\
\hline
\tlbreak{Charlotte}{InterruptRoute} & Interrupt binding, queued wake, unbinding,
    deferred drain & A queued wake cannot be delivered through a later route
    generation to the wrong owner. \\
\hline
\tlbreak{Charlotte}{Capability} & Unified per-address-space capability namespace
    and move rollback & Handles are fresh and authoritative; move revokes the
    source; rollback restores only the exact transaction; a closed namespace
    is empty. \\
\hline
\tlbreak{Charlotte}{Authorization} & Versioned subject/service policy,
    generation-fenced service bindings, decisions and capability issuance &
    Only authorized roles mutate policy or publication; decisions are bound to
    a subject, policy version, service generation and rights; redemption is
    single-use and cannot amplify authority. \\
\hline
\tlbreak{Charlotte}{DMA} & SMMU stream ownership, mapping, invalidation,
    quarantine, driver exit, pin reclamation & Unique stream/domain ownership;
    no cross-domain mapping; mapped memory stays pinned; failed hardware
    acknowledgement retains rather than prematurely frees authority. \\
\hline
\tlbreak{Charlotte}{ServiceLifecycle} & Signed artifact staging, loading, start,
    two-phase publication, lookup, fenced unregister, exit and teardown &
    Untrusted images never load; only activated generations resolve; stale
    activation/unregister cannot affect a replacement; teardown follows
    domain-wide reaping. \\
\hline
\tlbreak{Charlotte}{Raft} & Terms, durable votes, election, higher-term step-down,
    crash/restart & A voter records at most one durable vote per term and at
    most one leader can hold a term. \\
\hline
\tlbreak{Charlotte}{RaftLog} & Append, conflict repair, replication, commit,
    crash/restart & Log matching, committed-entry agreement, commit within the
    log, and leader completeness. \\
\hline
\tlbreak{Charlotte}{RaftMembership} & Voters/learners, joint consensus,
    finalization, decommissioning & Both voter majorities are required while
    joint; all proposed peers reach the joint fence before finalization;
    leadership and decommissioning follow active membership. \\
\hline
\tlbreak{Charlotte}{RaftJoin} & Selected-anchor admission before membership,
    JOIN fence, catch-up, crash/restart & A joiner accepts only the chosen
    anchor, cannot campaign, is promoted only after catch-up, and restores its
    durable admission fence after restart. \\
\hline
\tlbreak{Charlotte}{RaftSnapshot} & Chunk reception, durable installation,
    matching suffix, membership recovery & Stale snapshots cannot regress
    progress; the installed state and retained suffix agree; restart restores
    membership and application state before use. \\
\hline
\tlbreak{Charlotte}{RemoteCall} & Call identity, target generation, duplicate
    suppression, timeout uncertainty, result eviction & At most one execution
    while tracked; stale targets never execute; uncertain outcomes are not
    reported as success; deduplication evidence remains until transport
    settlement. \\
\hline
\tlbreak{Charlotte}{ReliableMessage} & Wire sessions across retry abandonment and
    unilateral service restart & Logical epochs have unique ordered wire
    identities and delayed sessions cannot move receive state backwards. \\
\hline
\end{longtable}
\normalsize

This decomposition also states what is \emph{not} silently being claimed. For
example, Ethernet fragmentation and retry timing are not part of the Raft
election proof; protobuf bytes and peer addresses are not part of joint
consensus; cryptographic correctness is not part of the service-lifecycle
trust-bit abstraction. Such omissions are recorded in the conformance map.

\section[Model-driven findings]{How model reasoning changes CharlotteOS}

The most valuable result is often a mismatch discovered while constructing the
model. The correct response depends on which side was wrong: sometimes the
manual or model overstated functionality; sometimes the code lacked a safety
fence.

\subsection{Authorization is controlled capability issuance}

The current name service has an optional reusable 64-bit access-key gate. It
does not identify a stable principal, express different rights per caller,
separate policy administration from service publication, or version a policy
decision. It is therefore not the production policy service suggested by the
architecture.

The authorization model makes the intended boundary explicit. The catalog
publishes a service generation and the maximum rights that may be delegated.
A separate logical policy role maps a stable principal and service identity to
an allowed-rights set and policy version. A decision is bound to the
kernel-authenticated subject, requested rights, current service generation and
policy version, and can be redeemed once for an attenuated connection. A
co-located name/policy process can perform decision and redemption as one
serialized transition; a later split service needs an unforgeable single-use
grant and must revalidate every binding at redemption.

The model also prevents an attractive but false claim. A policy update blocks
new issuance and invalidates unredeemed decisions, but cannot selectively
retract a direct connection already delegated by the kernel. Hard revocation
requires endpoint closure, a revocable proxy, or a future kernel lease object.
This prospective revocation contract is safer than documenting authority that
the implementation cannot enforce. The practical design is recorded in
\repofile|docs/architecture/authorization-policy.md|.

\subsection{Two-phase service publication}

An earlier lifecycle abstraction treated replacement publication as one
atomic swap. The implemented distributed catalog instead performs three
observable steps: commit an inactive prepared generation, publish the local
connection, and commit activation. Lookup intentionally excludes the inactive
generation, so replacement has a temporary unavailable interval.

The model was changed to describe that real contract rather than claim
continuous availability. It also exposed the code-level obligation that
generation allocation must fail closed. Distributed generation saturation and
local unchecked increment were replaced with checked allocation that returns
an error without mutating the catalog. A delayed unregister is accepted only
when both owner and generation still match, which preserves a replacement.

Formal safeguarding does not promise every desired availability property. Here it
prevents the documentation from claiming atomic replacement and protects the
generation safety that the implementation does provide.

\subsection{Reliable-message session identity}

The reliable-message service formerly initialized a wire session from the
service generation and incremented that number when abandoning a retry epoch.
Generation $N$ after one abandonment could therefore use the same identity as
generation $N+1$ before abandonment. A bounded retired-session window also
allowed a sufficiently old delayed SYN to replace newer receive state after it
fell out of the window.

The session model represents a session as the ordered pair
$(\mathit{serviceGeneration},\mathit{retryEpoch})$. The repaired wire encoding
packs those values into disjoint 32-bit fields, accepts a replacement session
only when the packed identity is well formed and strictly newer, and disables
sends rather than wrapping at exhaustion. Two unsafe configurations retain the
old flat-identity collision and receive-regression behaviors as required
counterexamples.

\subsection{Restart-safe Raft admission}

Dynamic admission originally kept \texttt{joining} and
\texttt{joining\_from} only in process memory. A node could select an existing
cluster as its authority, crash before membership became durable, restart from
its launch-time singleton configuration, and campaign. The gap represented
real CharlotteOS functionality, not merely a missing model action.

The repaired implementation persists a \texttt{JoinAdmission} record
containing the selected anchor and the snapshot index that existed before
admission. Construction restores the joining posture and disables elections.
When the joint membership commits, the node first writes a newer
membership-bearing snapshot and only then clears the durable admission record.
This ordering covers both crash edges:

\begin{itemize}
    \item before the membership snapshot is durable, restart retains the
        anchor fence and the node cannot campaign;
    \item after the snapshot is durable but before the fence-clear write,
        restart recognizes the strictly newer membership snapshot as completion
        evidence and safely clears the stale admission record.
\end{itemize}

Auto-join is disabled when the generic Raft service falls back to volatile
storage. That is an explicit safety-over-availability decision: without stable
state, a process restart cannot prove which cluster has authority. The safe
TLA+ specification includes crash and restart; an unsafe restart that forgets
admission must violate \texttt{RestartPreservesAdmission}.

\subsection{Retained regression patterns}

Other negative specifications preserve compact versions of repaired faults:

\small
\begin{longtable}{|p{0.27\textwidth}|p{0.31\textwidth}|p{0.32\textwidth}|}
\hline
\textbf{Unsafe action} & \textbf{Failure represented} & \textbf{Required violation} \\
\hline
\endfirsthead
\hline
\textbf{Unsafe action} & \textbf{Failure represented} & \textbf{Required violation} \\
\hline
\endhead
\tlbreak{Unsafe}{MessageWake} & Message readiness and endpoint-close observation
    share one observer queue. & \tlbreak{CloseSignalImplies}{Closed} \\
\hline
\tlbreak{Unsafe}{RemoteAbort} & A remote LP removes and reaps a context that may
    still execute. & \texttt{ReapOnlyOffCpu} \\
\hline
\texttt{UnsafeStaleClose} & A recycled numeric ASID is closed through an old
    generation handle. & \tlbreak{ReplacementSurvives}{StaleHandle} \\
\hline
\tlbreak{Unsafe}{Allocate} & A hardware ASID is reassigned before stale
    translations are invalidated. & \texttt{NoDirtyTagReuse} \\
\hline
\tlbreak{Drain}{Unsafe} & A deferred interrupt wake ignores its route generation.
    & \texttt{NoStaleWakeDelivery} \\
\hline
\tlbreak{UnsafeStale}{Unregister} & Cleanup from an old service generation removes
    its replacement. & \tlbreak{ReplacementSurvives}{StaleUnregister} \\
\hline
\tlbreak{UnsafeReplicate}{ToJoiner} & A non-selected peer supplies authority to a
    joining node. & \tlbreak{JoiningAcceptsOnly}{SelectedAnchor} \\
\hline
\tlbreak{UnsafeRestartForgets}{Admission} & Restart makes an incomplete joiner
    election-eligible. & \tlbreak{RestartPreserves}{Admission} \\
\hline
Flat session allocation & Retry and restart epochs alias one wire session. &
    \tlbreak{SessionIdentity}{Unique} \\
\hline
\texttt{AcceptDelayedSession} & An older delayed SYN regresses receive state. &
    \tlbreak{ReceiveSession}{Monotonic} \\
\hline
\texttt{UnsafeSetPolicy} & A non-administrator changes a policy rule. &
    \texttt{PolicyMutationAuthorized} \\
\hline
\tlbreak{UnsafeRedeemOther}{Principal} & A decision for one subject is redeemed
    for another. & \texttt{MintBoundToPrincipal} \\
\hline
\tlbreak{UnsafeRedeemStale}{Policy} & A decision survives a policy-version
    change. & \texttt{MintUsesCurrentPolicy} \\
\hline
\tlbreak{UnsafeRedeemStale}{Binding} & A decision for an old service generation
    reaches its replacement. & \texttt{MintTargetsCurrentBinding} \\
\hline
\tlbreak{UnsafeAmplify}{Rights} & Redemption mints rights absent from the
    decision. & \texttt{NoRightsAmplification} \\
\hline
\end{longtable}
\normalsize

The check suite expects all fifteen negative configurations to fail for the named
reason. If one unexpectedly satisfies its invariant, the regression harness
has lost sensitivity and the suite fails.

\section{Concrete conformance and atomicity}

TLA+ actions are atomic. Rust code usually crosses several structures, locks,
or persistent writes. The conformance review must therefore identify a
linearization point or represent the concrete intermediate states explicitly.

The mapping uses two correspondence classes:

\begin{description}
    \item[Direct] The abstract action follows one concrete transition while
        omitting data irrelevant to the invariant. For example, a successful
        fenced unregister directly checks owner and generation.
    \item[Abstract] Several concrete operations are collapsed. This is valid
        only when intermediate states are unobservable or separately guarded.
        For example, scheduler retirement spans multiple tables, so the
        concrete \texttt{RETIREMENTS\_IN\_FLIGHT} state prevents address-space
        teardown during that interval.
\end{description}

Persistent ordering deserves the same care. In Raft admission, ``snapshot
membership and clear join fence'' is one model transition, but two durable
writes in Rust. Safety depends on snapshot flush preceding fence removal. A
crash between those writes must map to a modeled state that remains safe.

The conformance document is not yet a machine-checked refinement proof. It is
a reviewable refinement \emph{obligation}: it names the functions and state
that must be re-examined when either side changes. Code review should reject a
model update that merely hides a new concrete intermediate state.

\section{Validation}

\subsection{Running TLC}

The repository wrapper runs the complete bounded suite:

\begin{verbatim}
docs/tla/check.sh /path/to/tla2tools.jar
\end{verbatim}

Alternatively, \texttt{TLA2TOOLS\_JAR} may name the tools archive. The script:

\begin{itemize}
    \item runs all eighteen safe configurations;
    \item enables action coverage and requires selected actions to have a
        positive successor count;
    \item rejects structural TLC warnings about variables, \texttt{EXCEPT}
        expressions, or incomplete successor states;
    \item runs fifteen negative configurations and requires each one to fail on
        the named invariant after exercising the named unsafe action; and
    \item writes checkpoints and traces into a temporary directory so generated
        state does not become part of the source tree.
\end{itemize}

Action coverage is important. An invariant can pass because a transition is
accidentally disabled. Requiring coverage does not prove that an action is
correct, but it detects this common vacuity.

At the 11 August 2026 implementation synchronization point, representative
fast configurations produced the following complete state graphs:

\begin{longtable}{|p{0.32\textwidth}|r|r|r|}
\hline
\textbf{Model} & \textbf{Generated} & \textbf{Distinct} & \textbf{Depth} \\
\hline
\texttt{CharlotteIPC} & 6,118,645 & 1,254,153 & 13 \\
\hline
\texttt{CharlotteCQ} & 1,688,159 & 321,967 & 12 \\
\hline
\texttt{CharlotteScheduler} & 1,710,157 & 373,860 & 22 \\
\hline
\tlbreak{Charlotte}{ServiceLifecycle} & 665,641 & 74,232 & 41 \\
\hline
\tlbreak{Charlotte}{Authorization} & 144,598 & 54,705 & 12 \\
\hline
\texttt{CharlotteRaftJoin} & 2,029 & 520 & 13 \\
\hline
\tlbreak{Charlotte}{ReliableMessage} & 30 & 21 & 5 \\
\hline
\end{longtable}

State counts are evidence about one configuration, not a quality score. A
small model may encode the decisive identity argument, while a large model may
contain substantial independent bookkeeping. The invariant, abstraction, and
coverage matter more than graph size.

\subsection{Implementation validation}

TLC is followed by validation proportional to the affected code. For the
August synchronization and admission repair this included:

\begin{itemize}
    \item host tests for \texttt{catten-graft}, including restart before durable
        membership and restart between membership persistence and fence clear;
    \item host tests for \texttt{charlotte-protocol-msg}, including disjoint
        retry/restart sessions and fail-closed exhaustion;
    \item AArch64 release builds and warning-denied Clippy for the changed
        \texttt{raft}, \texttt{relmsg}, and \texttt{ns} services;
    \item an AArch64 kernel check containing the updated persistent-storage join
        self-test manifest; and
    \item formatting, shell syntax, and diff-whitespace checks.
\end{itemize}

Boot and multi-guest tests remain necessary when a change affects the concrete
scheduler, device, storage, or network integration. A TLC result cannot show
that QEMU delivers an interrupt, that an object-store flush reaches the
emulated disk, or that a manifest starts the intended services.

\section{What a successful check means}

A successful safe configuration means that TLC found no reachable violation
of the declared invariants within that specification and its finite constants.
It does \emph{not} by itself establish:

\begin{itemize}
    \item a proof that arbitrary larger bounds preserve the result;
    \item a mechanically checked refinement from Rust to TLA+;
    \item liveness, fairness, bounded latency, or eventual leader election
        unless explicitly expressed as temporal properties;
    \item correctness of omitted bytes, parsers, cryptography, memory mappings,
        hardware programming, or transport framing;
    \item correctness under storage failure modes below the modeled durable
        write interface; or
    \item freedom from ordinary implementation bugs such as arithmetic,
        aliasing, or unsafe-memory errors outside the modeled projection.
\end{itemize}

Conversely, a gap between code and model is actionable information. It can
mean one of three things:

\begin{enumerate}
    \item \textbf{Functionality is unsafe.} The Raft admission restart gap was
        of this kind and required a CharlotteOS change.
    \item \textbf{The model overclaims.} Atomic service replacement was of this
        kind; the model was corrected to admit temporary unavailability.
    \item \textbf{The behavior is outside scope.} It must be documented rather
        than silently attributed to the checked invariants.
\end{enumerate}

Therefore a green TLC run is not a blanket ``formally verified'' label for
CharlotteOS.

\section{Workflow for changes}

Any change to ownership, generations, persistence, retry identity, lifecycle,
or distributed membership should be checked against this chapter's method.
A practical review checklist is:

\begin{enumerate}
    \item Search \repofile|docs/tla/CONFORMANCE.md| for the changed Rust
        function or state field.
    \item Decide whether the change adds an action, changes an enabling
        condition, moves a linearization point, or changes durable/volatile
        classification.
    \item Update the safe model and bounds. If repairing a discovered fault,
        retain the old behavior as an unsafe action and named negative
        configuration.
    \item Add the new action to the required-coverage list in
        \repofile|docs/tla/check.sh|.
    \item Run the entire TLC suite, not only the edited module, because shared
        documentation claims and layered Raft assumptions can drift together.
    \item Update the conformance map with exact Rust functions and ordering.
    \item Add implementation tests at the crash, retry, or stale-generation
        edges identified by the model.
    \item Run target validation and the relevant boot or multi-node suite.
    \item Record deliberate omissions and avoid converting a safety result into
        an unsupported liveness or availability claim.
\end{enumerate}

Used this way, TLA+ is part of the engineering control system for CharlotteOS.
It preserves the reasoning behind a safety fence, makes regressions executable,
and forces the model, manual, code, and tests to describe the same
functionality.

\endinput
